\begin{table}[!htbp]
\centering\zihao{-4}\setlength{\tabcolsep}{2.25pt}
\begin{minipage}{\textwidth}\songti\normalsize
\QTwoParagraph 这些方法分别通过张量交互、时序记忆、表征约束或辅助任务组织多模态信息，缺失增强方法还引入重构、表征对齐、元采样及潜在表示推断。比较同时考察完整输入预测能力、连续缺失后的表现和参数规模：前两项反映同一配置在不同输入条件下的结果，参数量补充表示模型复杂度。各方法接入统一特征与双任务接口，以下结果均限定在本文统一适配与评价设置下。
\end{minipage}\par\vspace{10bp}
\caption{公开方法与LTARP的完整输入验证集结果（三随机种子均值$\pm$样本标准差）}
\label{tab:q2_public_baselines}
\begin{tabularx}{\textwidth}{@{}>{\raggedright\arraybackslash}p{2.0cm}>{\centering\arraybackslash}p{2.5cm}>{\centering\arraybackslash}p{2.1cm}*{4}{>{\centering\arraybackslash}X}@{}}
\toprule
模型 & 训练方式 & 参数量 & 准确率$\uparrow$ & \shortstack{宏平均\\F1$\uparrow$} & MAE$\downarrow$ & \shortstack{相关系数\\$\uparrow$}\\
\midrule
TFN & \shortstack{完整输入\\训练} & 4,840,670 & \shortstack{$0.6099$\\$\pm0.0076$} & \shortstack{$0.5872$\\$\pm0.0083$} & \shortstack{$0.6336$\\$\pm0.0102$} & \shortstack{$0.6056$\\$\pm0.0098$}\\
LMF & \shortstack{完整输入\\训练} & 323,276 & \shortstack{$0.5925$\\$\pm0.0114$} & \shortstack{$0.5719$\\$\pm0.0090$} & \shortstack{$0.6682$\\$\pm0.0196$} & \shortstack{$0.5807$\\$\pm0.0188$}\\
MFN & \shortstack{完整输入\\训练} & 261,764 & \shortstack{$0.5966$\\$\pm0.0080$} & \shortstack{$0.5824$\\$\pm0.0062$} & \shortstack{$0.6418$\\$\pm0.0124$} & \shortstack{$0.5908$\\$\pm0.0125$}\\
MulT & \shortstack{完整输入\\训练} & 624,094 & \shortstack{$0.6172$\\$\pm0.0148$} & \shortstack{$0.5891$\\$\pm0.0236$} & \shortstack{$0.6221$\\$\pm0.0260$} & \shortstack{$0.6369$\\$\pm0.0115$}\\
MISA & \shortstack{完整输入\\训练} & 1,118,852 & \shortstack{$0.6131$\\$\pm0.0083$} & \shortstack{$0.6010$\\$\pm0.0060$} & \shortstack{$0.9175$\\$\pm0.1005$} & \shortstack{$0.6269$\\$\pm0.0087$}\\
Self-MM & \shortstack{完整输入\\训练} & 82,119 & \shortstack{$0.6355$\\$\pm0.0052$} & \shortstack{$0.6107$\\$\pm0.0084$} & \shortstack{$0.6099$\\$\pm0.0066$} & \shortstack{$0.6423$\\$\pm0.0042$}\\
MMIM & \shortstack{完整输入\\训练} & 163,204 & \shortstack{$0.6360$\\$\pm0.0048$} & \shortstack{$0.6110$\\$\pm0.0007$} & \shortstack{$0.6102$\\$\pm0.0098$} & \shortstack{$0.6375$\\$\pm0.0054$}\\
\addlinespace\midrule
TFR-Net & \shortstack{缺失增强\\训练} & 415,620 & \shortstack{$0.6374$\\$\pm0.0071$} & \shortstack{$0.6230$\\$\pm0.0080$} & \shortstack{$0.6145$\\$\pm0.0118$} & \shortstack{$0.6267$\\$\pm0.0142$}\\
MissModal & \shortstack{缺失增强\\训练} & 163,204 & \shortstack{$0.6305$\\$\pm0.0071$} & \shortstack{$0.6108$\\$\pm0.0069$} & \shortstack{$0.6084$\\$\pm0.0113$} & \shortstack{$0.6431$\\$\pm0.0006$}\\
M3S & \shortstack{缺失增强\\训练} & 323,276 & \shortstack{$0.6035$\\$\pm0.0063$} & \shortstack{$0.5775$\\$\pm0.0169$} & \shortstack{$0.6706$\\$\pm0.0173$} & \shortstack{$0.5793$\\$\pm0.0171$}\\
MMIN & \shortstack{缺失增强\\训练} & 292,164 & \shortstack{$0.6383$\\$\pm0.0062$} & \shortstack{$0.6117$\\$\pm0.0079$} & \shortstack{$0.6279$\\$\pm0.0247$} & \shortstack{$0.6401$\\$\pm0.0027$}\\
\addlinespace\midrule
LTARP & \shortstack{完整输入\\训练} & 164,343 & \shortstack{$0.6415$\\$\pm0.0050$} & \shortstack{$0.6243$\\$\pm0.0026$} & \shortstack{$0.6052$\\$\pm0.0063$} & \shortstack{$0.6475$\\$\pm0.0024$}\\
\bottomrule
\end{tabularx}
\par\vspace{10bp}\begin{minipage}{\textwidth}\songti\normalsize\setlength{\parindent}{2em}
\QTwoParagraph 在本文统一适配与评价设置下，LTARP在完整输入的四项指标上均取得表~\ref{tab:q2_public_baselines}中最优的三随机种子均值。部分轻量方法及缺失增强方法的表现接近LTARP，而更复杂的融合结构未呈现同等的平均优势。Self-MM和MMIM借助辅助任务或信息约束，在较小规模下也保留了较强预测信息；TFR-Net的宏平均 F1接近LTARP，但回归指标有所差异。因此，分类均衡性和回归结果需要结合观察。
\end{minipage}
\end{table}